% Tank radius 1, bottom y=0, rim y=2, outlet y=3.
\begin{tikzpicture}[x=1.35cm,y=1.1cm,line width=.7pt,
  every node/.style={font=\small,inner sep=2pt}]
% A transparent cylindrical tank, including its curved bottom.
\fill[figureblue!5] (-1,2)--(-1,0)
  arc[start angle=180,end angle=360,x radius=1,y radius=.22]
  --(1,2) arc[start angle=0,end angle=180,x radius=1,y radius=.22];
\draw[gray,dashed] (-1,0)
  arc[start angle=180,end angle=0,x radius=1,y radius=.22];
% The selected layer is a thin circular disk, not a rectangular strip.
\fill[figureblue!28] (-1,.95)
  arc[start angle=180,end angle=360,x radius=1,y radius=.22]
  --(1,.85) arc[start angle=0,end angle=-180,x radius=1,y radius=.22]--cycle;
\filldraw[fill=figureblue!15,draw=figureaccent]
  (0,.95) ellipse[x radius=1,y radius=.22];
\draw[figureaccent,dashed] (-1,.85)
  arc[start angle=180,end angle=0,x radius=1,y radius=.22];
\draw[figureaccent] (-1,.95)--(-1,.85)
  arc[start angle=180,end angle=360,x radius=1,y radius=.22]--(1,.95);
\draw (-1,2)--(-1,0)
  arc[start angle=180,end angle=360,x radius=1,y radius=.22]--(1,2);
\draw (0,2) ellipse[x radius=1,y radius=.22];
% Radius of the circular tank opening.
\fill (0,2) circle[radius=.025];
\draw[->] (0,2)--(1,2);
\node[fill=figureblue!5,inner sep=1pt] at (.5,2) {$1$};
% Height scale and the level of the outlet.
\draw[gray,dashed] (-1.7,3)--(2.95,3);
\draw[gray,dashed] (-1.7,2)--(-1.05,2)
  (-1.7,.9)--(-1.05,.9) (-1.7,0)--(-1.05,0);
\draw[->] (-1.75,-.35)--(-1.75,3.4);
\node[above] at (-1.75,3.4) {$y$};
\foreach \h in {0,2,3}
  \draw (-1.81,\h)--(-1.69,\h) node[left=4pt] {$\h$};
\draw (-1.81,.9)--(-1.69,.9) node[left=4pt] {$y$};
% A hollow outlet pipe: its lip lies at y=3.
\draw[gray!80,line width=3.5pt,line join=round]
  (1,2)--(1.65,2)--(1.65,3)--(2.35,3);
\draw[white,line width=2pt,line join=round]
  (1,2)--(1.65,2)--(1.65,3)--(2.35,3);
\draw[->,figureaccent] (2.17,3)--(2.55,3);
\node[above=4pt] at (2.05,3) {outlet};
% External arrowheads keep the thin layer's thickness readable.
\draw[gray,dashed] (1,.85)--(1.4,.85) (1,.95)--(1.4,.95);
\draw[->] (1.22,.5)--(1.22,.85);
\draw[->] (1.22,1.3)--(1.22,.95);
\node[right] at (1.22,1.2) {$dy$};
% The lift is measured from the selected layer up to the outlet.
\draw[gray,dashed] (1.05,.9)--(2.95,.9);
\draw[<->] (2.8,.9)--(2.8,3);
\node[right=3pt] at (2.8,1.95) {$3-y$};
\end{tikzpicture}
